
\documentclass{article}
\usepackage{colortbl}
\usepackage{makecell}
\usepackage{multirow}
\usepackage{supertabular}

\begin{document}

\newcounter{utterance}

\centering \large Interaction Transcript for game `function\_detective', experiment `string\_passive\_labeled\_pairs', episode 159 with mistral{-}large{-}3{-}azure.
\vspace{24pt}

{ \footnotesize  \setcounter{utterance}{1}
\setlength{\tabcolsep}{0pt}
\begin{supertabular}{c@{$\;$}|p{.15\linewidth}@{}p{.15\linewidth}p{.15\linewidth}p{.15\linewidth}p{.15\linewidth}p{.15\linewidth}}
    \# & \multicolumn{2}{c}{Player} && \multicolumn{2}{c}{Game Master} \\
    \hline

    \theutterance \stepcounter{utterance}  
    & & & \multicolumn{4}{p{0.6\linewidth}}{
        \cellcolor[rgb]{0.9,0.9,0.9}{
            \makecell[{{p{\linewidth}}}]{
                \texttt{\tiny{[P1$\langle$GM]}}
                \texttt{I have a hidden Python function for you to figure out.} \\
\\ 
\texttt{Your job is to infer what it does and then submit a correct Python implementation.} \\
\\ 
\texttt{Function signature:} \\
\texttt{(text: str) {-}> str} \\
\\ 
\texttt{You have 15 turns total.} \\
\\ 
\texttt{I reveal candidate input{-}output pairs one at a time and tell you whether each one is correct.} \\
\\ 
\texttt{You must make your best SOLVE attempt no later than the second{-}to{-}last turn.} \\
\texttt{Do not wait until the very end.} \\
\\ 
\texttt{You can use these actions:} \\
\texttt{{-} JSON format described at the end of the prompt} \\
\texttt{{-} SOLVE: ```python ... ```} \\
\\ 
\texttt{Do not add import statements in your guess.} \\
\texttt{Anything you might need is already available, so just write the function body in hypothesis and return the answer.} \\
\\ 
\texttt{Preloaded examples:} \\
\texttt{1. (('   ', '20 20 20'), True)} \\
\\ 
\texttt{Scoring:} \\
\texttt{{-} correctness is required for any non{-}zero score} \\
\texttt{{-} in interactive modes, solving with fewer information requests is better} \\
\\ 
\texttt{Input formatting reminders:} \\
\texttt{{-} strings must be quoted} \\
\texttt{{-} booleans must be True/False} \\
\texttt{{-} use commas between multiple values} \\
\texttt{{-} do not write any extra action tags in the same message} \\
\\ 
\texttt{Format example:} \\
\texttt{You write:} \\
\texttt{\{'current\_mode': <probing/confirming/don't know; choose whichever best describes your current state in the task>, 'input\_rationale': <brief description of why you chose this hypothesis>, 'hypothesis': <either the exact string "unknown" if you have no guess yet OR your current best guess as runnable Python matching the given signature, with a one{-}line docstring stating what it does. Never write placeholder code, if you cannot commit to a guess, use "unknown". Match the given signature.} \\
\texttt{e.g.} \\
\texttt{def f(x):} \\
\texttt{	"""Returns the absolute value of x."""} \\
\texttt{	return x if x >= 0 else {-}x} \\
\\ 
\texttt{Use "unknown" if you have no guess.>, 'confidence\_score': <score between 1{-}100. 100 means highest confidence in your hypothesis>\}} \\
\\ 
\texttt{Output only one valid json per turn.} \\
            }
        }
    }
     \\ \\

    \theutterance \stepcounter{utterance}  
    & \multicolumn{4}{p{0.6\linewidth}}{
        \cellcolor[rgb]{0.561,0.737,0.561}{
            \makecell[{{p{\linewidth}}}]{
                \texttt{\tiny{[P1$\rangle$GM]}}
                \texttt{```json} \\
\texttt{\{} \\
\texttt{  "current\_mode": "probing",} \\
\texttt{  "input\_rationale": "The preloaded example shows that the input string '   ' (three spaces) is transformed into '20 20 20'. This suggests that the function might be converting each character in the input string into its ASCII code, as the ASCII code for a space is 32, but 20 is the hexadecimal representation of 32. Alternatively, it could be converting to URL encoding, where a space is \%20, but the output is '20 20 20' without the '\%'. The most plausible hypothesis is that the function converts each character to its hexadecimal ASCII value, separated by spaces.",} \\
\texttt{  "hypothesis": "def f(text: str) {-}> str:\textbackslash{}n    \textbackslash{}"\textbackslash{}"\textbackslash{}"Converts each character in the input string to its hexadecimal ASCII value, separated by spaces.\textbackslash{}"\textbackslash{}"\textbackslash{}"\textbackslash{}n    return ' '.join(f'\{ord(c):02x\}' for c in text)",} \\
\texttt{  "confidence\_score": 70} \\
\texttt{\}} \\
\texttt{```} \\
            }
        }
    }
    & & \\ \\

    \theutterance \stepcounter{utterance}  
    & & & \multicolumn{4}{p{0.6\linewidth}}{
        \cellcolor[rgb]{0.9,0.9,0.9}{
            \makecell[{{p{\linewidth}}}]{
                \texttt{\tiny{[P1$\langle$GM]}}
                \texttt{OUTPUT: (('ab', '61 62\_alt'), False)} \\
            }
        }
    }
     \\ \\

    \theutterance \stepcounter{utterance}  
    & \multicolumn{4}{p{0.6\linewidth}}{
        \cellcolor[rgb]{0.561,0.737,0.561}{
            \makecell[{{p{\linewidth}}}]{
                \texttt{\tiny{[P1$\rangle$GM]}}
                \texttt{```json} \\
\texttt{\{} \\
\texttt{  "current\_mode": "probing",} \\
\texttt{  "input\_rationale": "The previous hypothesis was incorrect for the input 'ab', which should output '61 62\_alt' instead of '61 62'. This suggests that the function is not simply converting each character to its hexadecimal ASCII value. The '\_alt' suffix indicates that there might be a conditional transformation, possibly distinguishing between certain types of characters (e.g., letters vs. non{-}letters, or uppercase vs. lowercase). Alternatively, it could be related to URL encoding or another encoding scheme where certain characters are treated specially. I will revise the hypothesis to account for this '\_alt' suffix.",} \\
\texttt{  "hypothesis": "def f(text: str) {-}> str:\textbackslash{}n    \textbackslash{}"\textbackslash{}"\textbackslash{}"Converts each character to its hexadecimal ASCII value, appending '\_alt' for alphabetic characters.\textbackslash{}"\textbackslash{}"\textbackslash{}"\textbackslash{}n    result = {[}{]}\textbackslash{}n    for c in text:\textbackslash{}n        hex\_val = f'\{ord(c):02x\}'\textbackslash{}n        if c.isalpha():\textbackslash{}n            hex\_val += '\_alt'\textbackslash{}n        result.append(hex\_val)\textbackslash{}n    return ' '.join(result)",} \\
\texttt{  "confidence\_score": 60} \\
\texttt{\}} \\
\texttt{```} \\
            }
        }
    }
    & & \\ \\

    \theutterance \stepcounter{utterance}  
    & & & \multicolumn{4}{p{0.6\linewidth}}{
        \cellcolor[rgb]{0.9,0.9,0.9}{
            \makecell[{{p{\linewidth}}}]{
                \texttt{\tiny{[P1$\langle$GM]}}
                \texttt{OUTPUT: (('a{-}b', '61 2d 62'), True)} \\
            }
        }
    }
     \\ \\

    \theutterance \stepcounter{utterance}  
    & \multicolumn{4}{p{0.6\linewidth}}{
        \cellcolor[rgb]{0.561,0.737,0.561}{
            \makecell[{{p{\linewidth}}}]{
                \texttt{\tiny{[P1$\rangle$GM]}}
                \texttt{```json} \\
\texttt{\{} \\
\texttt{  "current\_mode": "confirming",} \\
\texttt{  "input\_rationale": "The previous hypothesis was incorrect because the input 'a{-}b' outputs '61 2d 62' without any '\_alt' suffix, which contradicts the idea that alphabetic characters always get '\_alt'. However, the preloaded example ('   ' {-}> '20 20 20') and this new example ('a{-}b' {-}> '61 2d 62') suggest that the function is simply converting each character to its hexadecimal ASCII value, with no additional suffixes. The '\_alt' in the earlier incorrect output might have been a red herring or a special case for certain inputs. I will revert to the original hypothesis of pure hexadecimal conversion.",} \\
\texttt{  "hypothesis": "def f(text: str) {-}> str:\textbackslash{}n    \textbackslash{}"\textbackslash{}"\textbackslash{}"Converts each character in the input string to its hexadecimal ASCII value, separated by spaces.\textbackslash{}"\textbackslash{}"\textbackslash{}"\textbackslash{}n    return ' '.join(f'\{ord(c):02x\}' for c in text)",} \\
\texttt{  "confidence\_score": 90} \\
\texttt{\}} \\
\texttt{```} \\
            }
        }
    }
    & & \\ \\

    \theutterance \stepcounter{utterance}  
    & & & \multicolumn{4}{p{0.6\linewidth}}{
        \cellcolor[rgb]{0.9,0.9,0.9}{
            \makecell[{{p{\linewidth}}}]{
                \texttt{\tiny{[P1$\langle$GM]}}
                \texttt{OUTPUT: (('aA', '61 41\_alt'), False)} \\
            }
        }
    }
     \\ \\

    \theutterance \stepcounter{utterance}  
    & \multicolumn{4}{p{0.6\linewidth}}{
        \cellcolor[rgb]{0.561,0.737,0.561}{
            \makecell[{{p{\linewidth}}}]{
                \texttt{\tiny{[P1$\rangle$GM]}}
                \texttt{```json} \\
\texttt{\{} \\
\texttt{  "current\_mode": "probing",} \\
\texttt{  "input\_rationale": "The output for 'aA' is '61 41\_alt', which shows that lowercase letters are converted to their hex ASCII values without a suffix, while uppercase letters get the '\_alt' suffix. This suggests the function distinguishes between lowercase and uppercase letters, appending '\_alt' only for uppercase letters. Non{-}alphabetic characters (like spaces or hyphens) do not get the suffix. This aligns with the previous examples: spaces ('20 20 20') and hyphens ('2d') had no suffix, while uppercase 'A' does.",} \\
\texttt{  "hypothesis": "def f(text: str) {-}> str:\textbackslash{}n    \textbackslash{}"\textbackslash{}"\textbackslash{}"Converts each character to its hexadecimal ASCII value, appending '\_alt' for uppercase letters.\textbackslash{}"\textbackslash{}"\textbackslash{}"\textbackslash{}n    result = {[}{]}\textbackslash{}n    for c in text:\textbackslash{}n        hex\_val = f'\{ord(c):02x\}'\textbackslash{}n        if c.isupper():\textbackslash{}n            hex\_val += '\_alt'\textbackslash{}n        result.append(hex\_val)\textbackslash{}n    return ' '.join(result)",} \\
\texttt{  "confidence\_score": 85} \\
\texttt{\}} \\
\texttt{```} \\
            }
        }
    }
    & & \\ \\

    \theutterance \stepcounter{utterance}  
    & & & \multicolumn{4}{p{0.6\linewidth}}{
        \cellcolor[rgb]{0.9,0.9,0.9}{
            \makecell[{{p{\linewidth}}}]{
                \texttt{\tiny{[P1$\langle$GM]}}
                \texttt{OUTPUT: (('123', '31 32 33'), True)} \\
            }
        }
    }
     \\ \\

    \theutterance \stepcounter{utterance}  
    & \multicolumn{4}{p{0.6\linewidth}}{
        \cellcolor[rgb]{0.561,0.737,0.561}{
            \makecell[{{p{\linewidth}}}]{
                \texttt{\tiny{[P1$\rangle$GM]}}
                \texttt{```json} \\
\texttt{\{} \\
\texttt{  "current\_mode": "confirming",} \\
\texttt{  "input\_rationale": "The output for '123' is '31 32 33', which matches the hypothesis that the function converts each character to its hexadecimal ASCII value without any suffix for non{-}uppercase letters. This confirms that digits and lowercase letters are treated the same way (no suffix), while uppercase letters get the '\_alt' suffix. The preloaded example ('   ' {-}> '20 20 20'), 'a{-}b' {-}> '61 2d 62', and '123' {-}> '31 32 33' all align with this behavior.",} \\
\texttt{  "hypothesis": "def f(text: str) {-}> str:\textbackslash{}n    \textbackslash{}"\textbackslash{}"\textbackslash{}"Converts each character to its hexadecimal ASCII value, appending '\_alt' for uppercase letters.\textbackslash{}"\textbackslash{}"\textbackslash{}"\textbackslash{}n    result = {[}{]}\textbackslash{}n    for c in text:\textbackslash{}n        hex\_val = f'\{ord(c):02x\}'\textbackslash{}n        if c.isupper():\textbackslash{}n            hex\_val += '\_alt'\textbackslash{}n        result.append(hex\_val)\textbackslash{}n    return ' '.join(result)",} \\
\texttt{  "confidence\_score": 95} \\
\texttt{\}} \\
\texttt{```} \\
            }
        }
    }
    & & \\ \\

    \theutterance \stepcounter{utterance}  
    & & & \multicolumn{4}{p{0.6\linewidth}}{
        \cellcolor[rgb]{0.9,0.9,0.9}{
            \makecell[{{p{\linewidth}}}]{
                \texttt{\tiny{[P1$\langle$GM]}}
                \texttt{OUTPUT: (('hello world', '68 65 6c 6c 6f 20 77 6f 72 6c 64\_alt'), False)} \\
            }
        }
    }
     \\ \\

    \theutterance \stepcounter{utterance}  
    & \multicolumn{4}{p{0.6\linewidth}}{
        \cellcolor[rgb]{0.561,0.737,0.561}{
            \makecell[{{p{\linewidth}}}]{
                \texttt{\tiny{[P1$\rangle$GM]}}
                \texttt{```json} \\
\texttt{\{} \\
\texttt{  "current\_mode": "probing",} \\
\texttt{  "input\_rationale": "The output for 'hello world' is '68 65 6c 6c 6f 20 77 6f 72 6c 64\_alt'. This shows that only the last character 'd' (which is lowercase) incorrectly got the '\_alt' suffix, while the uppercase 'W' in 'world' did not. This contradicts the previous hypothesis that only uppercase letters get '\_alt'. It seems the rule is more nuanced: perhaps only the last character gets '\_alt' if it is alphabetic (regardless of case), or there is another condition. I will adjust the hypothesis to append '\_alt' only to the last character if it is a letter.",} \\
\texttt{  "hypothesis": "def f(text: str) {-}> str:\textbackslash{}n    \textbackslash{}"\textbackslash{}"\textbackslash{}"Converts each character to its hexadecimal ASCII value, appending '\_alt' to the last character if it is a letter.\textbackslash{}"\textbackslash{}"\textbackslash{}"\textbackslash{}n    if not text:\textbackslash{}n        return ''\textbackslash{}n    result = {[}{]}\textbackslash{}n    for i, c in enumerate(text):\textbackslash{}n        hex\_val = f'\{ord(c):02x\}'\textbackslash{}n        if i == len(text) {-} 1 and c.isalpha():\textbackslash{}n            hex\_val += '\_alt'\textbackslash{}n        result.append(hex\_val)\textbackslash{}n} \\
            }
        }
    }
    & & \\ \\

    \theutterance \stepcounter{utterance}  
    & & & \multicolumn{4}{p{0.6\linewidth}}{
        \cellcolor[rgb]{0.9,0.9,0.9}{
            \makecell[{{p{\linewidth}}}]{
                \texttt{\tiny{[P1$\langle$GM]}}
                \texttt{Game Violation: ParseError: Response is not valid JSON.} \\
\texttt{Output ONLY:} \\
\texttt{One JSON with confidence\_score, current\_mode, input\_rationale, and hypothesis keys ...} \\
\texttt{OR} \\
\texttt{SOLVE: ```python ... ```} \\
            }
        }
    }
     \\ \\

    \theutterance \stepcounter{utterance}  
    & \multicolumn{4}{p{0.6\linewidth}}{
        \cellcolor[rgb]{0.561,0.737,0.561}{
            \makecell[{{p{\linewidth}}}]{
                \texttt{\tiny{[P1$\rangle$GM]}}
                \texttt{```json} \\
\texttt{\{} \\
\texttt{  "current\_mode": "probing",} \\
\texttt{  "input\_rationale": "The output for 'hello world' shows that only the last character ('d') gets the '\_alt' suffix, regardless of case. This suggests the function appends '\_alt' to the last character if it is alphabetic. All other characters (including uppercase letters) are converted to their hexadecimal ASCII values without suffixes. I will finalize this hypothesis.",} \\
\texttt{  "hypothesis": "def f(text: str) {-}> str:\textbackslash{}n    \textbackslash{}"\textbackslash{}"\textbackslash{}"Converts each character to its hexadecimal ASCII value, appending '\_alt' to the last character if it is a letter.\textbackslash{}"\textbackslash{}"\textbackslash{}"\textbackslash{}n    if not text:\textbackslash{}n        return ''\textbackslash{}n    result = {[}{]}\textbackslash{}n    for i, c in enumerate(text):\textbackslash{}n        hex\_val = f'\{ord(c):02x\}'\textbackslash{}n        if i == len(text) {-} 1 and c.isalpha():\textbackslash{}n            hex\_val += '\_alt'\textbackslash{}n        result.append(hex\_val)\textbackslash{}n    return ' '.join(result)",} \\
\texttt{  "confidence\_score": 95} \\
\texttt{\}} \\
\texttt{```} \\
            }
        }
    }
    & & \\ \\

    \theutterance \stepcounter{utterance}  
    & & & \multicolumn{4}{p{0.6\linewidth}}{
        \cellcolor[rgb]{0.9,0.9,0.9}{
            \makecell[{{p{\linewidth}}}]{
                \texttt{\tiny{[P1$\langle$GM]}}
                \texttt{OUTPUT: (('', ''), True)} \\
            }
        }
    }
     \\ \\

    \theutterance \stepcounter{utterance}  
    & \multicolumn{4}{p{0.6\linewidth}}{
        \cellcolor[rgb]{0.561,0.737,0.561}{
            \makecell[{{p{\linewidth}}}]{
                \texttt{\tiny{[P1$\rangle$GM]}}
                \texttt{```json} \\
\texttt{\{} \\
\texttt{  "current\_mode": "confirming",} \\
\texttt{  "input\_rationale": "The empty string case confirms the hypothesis handles empty input correctly. All previous test cases now align with the rule: convert each character to its hexadecimal ASCII value, and append '\_alt' to the last character if it's alphabetic. This explains all observed behavior: spaces ('20 20 20'), 'a{-}b' ('61 2d 62'), '123' ('31 32 33'), 'aA' ('61 41\_alt'), and 'hello world' ('68 65 6c 6c 6f 20 77 6f 72 6c 64\_alt'). I will now submit the final solution.",} \\
\texttt{  "hypothesis": "SOLVE: ```python\textbackslash{}ndef f(text: str) {-}> str:\textbackslash{}n    \textbackslash{}"\textbackslash{}"\textbackslash{}"Converts each character to its hexadecimal ASCII value, appending '\_alt' to the last character if it is a letter.\textbackslash{}"\textbackslash{}"\textbackslash{}"\textbackslash{}n    if not text:\textbackslash{}n        return ''\textbackslash{}n    result = {[}{]}\textbackslash{}n    for i, c in enumerate(text):\textbackslash{}n        hex\_val = f'\{ord(c):02x\}'\textbackslash{}n        if i == len(text) {-} 1 and c.isalpha():\textbackslash{}n            hex\_val +=} \\
            }
        }
    }
    & & \\ \\

    \theutterance \stepcounter{utterance}  
    & & & \multicolumn{4}{p{0.6\linewidth}}{
        \cellcolor[rgb]{0.9,0.9,0.9}{
            \makecell[{{p{\linewidth}}}]{
                \texttt{\tiny{[P1$\langle$GM]}}
                \texttt{Game Violation: ParseError: Response is not valid JSON.} \\
\texttt{Output ONLY:} \\
\texttt{One JSON with confidence\_score, current\_mode, input\_rationale, and hypothesis keys ...} \\
\texttt{OR} \\
\texttt{SOLVE: ```python ... ```} \\
            }
        }
    }
     \\ \\

    \theutterance \stepcounter{utterance}  
    & \multicolumn{4}{p{0.6\linewidth}}{
        \cellcolor[rgb]{0.561,0.737,0.561}{
            \makecell[{{p{\linewidth}}}]{
                \texttt{\tiny{[P1$\rangle$GM]}}
                \texttt{```python} \\
\texttt{SOLVE:} \\
\texttt{def f(text: str) {-}> str:} \\
\texttt{    """Converts each character to its hexadecimal ASCII value, appending '\_alt' to the last character if it is a letter."""} \\
\texttt{    if not text:} \\
\texttt{        return ''} \\
\texttt{    result = {[}{]}} \\
\texttt{    for i, c in enumerate(text):} \\
\texttt{        hex\_val = f'\{ord(c):02x\}'} \\
\texttt{        if i == len(text) {-} 1 and c.isalpha():} \\
\texttt{            hex\_val += '\_alt'} \\
\texttt{        result.append(hex\_val)} \\
\texttt{    return ' '.join(result)} \\
\texttt{```} \\
            }
        }
    }
    & & \\ \\

    \theutterance \stepcounter{utterance}  
    & & & \multicolumn{4}{p{0.6\linewidth}}{
        \cellcolor[rgb]{0.9,0.9,0.9}{
            \makecell[{{p{\linewidth}}}]{
                \texttt{\tiny{[P1$\langle$GM]}}
                \texttt{Game Violation: ParseError: Response is not valid JSON.} \\
\texttt{Output ONLY:} \\
\texttt{One JSON with confidence\_score, current\_mode, input\_rationale, and hypothesis keys ...} \\
\texttt{OR} \\
\texttt{SOLVE: ```python ... ```} \\
            }
        }
    }
     \\ \\

    \theutterance \stepcounter{utterance}  
    & \multicolumn{4}{p{0.6\linewidth}}{
        \cellcolor[rgb]{0.561,0.737,0.561}{
            \makecell[{{p{\linewidth}}}]{
                \texttt{\tiny{[P1$\rangle$GM]}}
                \texttt{SOLVE:} \\
\texttt{```python} \\
\texttt{def f(text: str) {-}> str:} \\
\texttt{    """Converts each character to its hexadecimal ASCII value, appending '\_alt' to the last character if it is a letter."""} \\
\texttt{    if not text:} \\
\texttt{        return ''} \\
\texttt{    parts = {[}{]}} \\
\texttt{    for i, c in enumerate(text):} \\
\texttt{        hex\_val = f"\{ord(c):02x\}"} \\
\texttt{        if i == len(text) {-} 1 and c.isalpha():} \\
\texttt{            hex\_val += "\_alt"} \\
\texttt{        parts.append(hex\_val)} \\
\texttt{    return " ".join(parts)} \\
\texttt{```} \\
            }
        }
    }
    & & \\ \\

    \theutterance \stepcounter{utterance}  
    & & & \multicolumn{2}{p{0.3\linewidth}}{
        \cellcolor[rgb]{0.9,0.9,0.9}{
            \makecell[{{p{\linewidth}}}]{
                \texttt{\tiny{[GM$|$GM]}}
                \texttt{Game Verdict: LOSS} \\
            }
        }
    }
    & & \\ \\

    \theutterance \stepcounter{utterance}  
    & & & \multicolumn{2}{p{0.3\linewidth}}{
        \cellcolor[rgb]{0.9,0.9,0.9}{
            \makecell[{{p{\linewidth}}}]{
                \texttt{\tiny{[GM$|$GM]}}
                \texttt{Incorrect.} \\
\texttt{Failed Cases:} \\
\texttt{Input: {[}'ab'{]} | Expected: 61 62 | Got: 61 62\_alt} \\
\texttt{Input: {[}'a{-}b'{]} | Expected: 61 2d 62 | Got: 61 2d 62\_alt} \\
\texttt{Input: {[}'aA'{]} | Expected: 61 41 | Got: 61 41\_alt} \\
            }
        }
    }
    & & \\ \\

    \theutterance \stepcounter{utterance}  
    & & & \multicolumn{2}{p{0.3\linewidth}}{
        \cellcolor[rgb]{0.9,0.9,0.9}{
            \makecell[{{p{\linewidth}}}]{
                \texttt{\tiny{[GM$|$GM]}}
                \texttt{def hex\_code\_points(text: str) {-}> str:} \\
\texttt{    return ' '.join(format(ord(ch), 'x') for ch in text)} \\
            }
        }
    }
    & & \\ \\

    \theutterance \stepcounter{utterance}  
    & & & \multicolumn{2}{p{0.3\linewidth}}{
        \cellcolor[rgb]{0.9,0.9,0.9}{
            \makecell[{{p{\linewidth}}}]{
                \texttt{\tiny{[GM$|$GM]}}
                \texttt{def f(text: str) {-}> str:} \\
\texttt{    """Converts each character to its hexadecimal ASCII value, appending '\_alt' to the last character if it is a letter."""} \\
\texttt{    if not text:} \\
\texttt{        return ''} \\
\texttt{    parts = {[}{]}} \\
\texttt{    for i, c in enumerate(text):} \\
\texttt{        hex\_val = f"\{ord(c):02x\}"} \\
\texttt{        if i == len(text) {-} 1 and c.isalpha():} \\
\texttt{            hex\_val += "\_alt"} \\
\texttt{        parts.append(hex\_val)} \\
\texttt{    return " ".join(parts)} \\
            }
        }
    }
    & & \\ \\

    \theutterance \stepcounter{utterance}  
    & & & \multicolumn{2}{p{0.3\linewidth}}{
        \cellcolor[rgb]{0.9,0.9,0.9}{
            \makecell[{{p{\linewidth}}}]{
                \texttt{\tiny{[GM$|$GM]}}
                \texttt{Mode: passive\_labeled\_pairs} \\
\texttt{Accuracy: 0.210} \\
\texttt{Information requests: 6/15} \\
\texttt{Revealed passive examples: 7} \\
\texttt{Efficiency: 0.667} \\
\texttt{Slack: 9} \\
\texttt{Observed{-}pair consistency: 0.750} \\
\texttt{Membership balance: 4 true / 3 false} \\
\texttt{Parse errors: 3} \\
\texttt{Runtime errors: 0} \\
            }
        }
    }
    & & \\ \\

\end{supertabular}
}

\end{document}
